
%%% Local Variables:
%%% mode: Xelatex
%%% TeX-master: t
%%% End:


%%% \ctitle{面向医疗多智能体系统的对抗性攻击与自动化安全评估框架研究}
\ctitle{医疗大语言模型安全评估框架与对抗鲁棒性研究}

\fenleihao{} \xuehao{M202373434} 
\schoolcode{10487} \miji{公开} %内部/秘密/机密/绝密
\csubjectname{控制科学与工程} \cauthorname{刘子恒}
\csupervisorname{夏天} \csupervisortitle{教授}
\defencedate{2026~年~4~月~29~日} \grantdate{}
\chair{}%
\firstreviewer{} \secondreviewer{} \thirdreviewer{}


\edegree{the Master Degree in Engineering} %（工学硕士）
%\edegree{the Master Degree in Science}     %（理学硕士）
%\edegree{the Professional Master Degree}   %(专业学位)
\etitle{Research on Safety Evaluation Framework and Adversarial Robustness of Medical Large Language Models}
\esubject{Control Science and Engineering}
\eauthor{LIU Ziheng}
\esupervisor{Prof. XIA Tian}

%定义中英文摘要和关键字
\cabstract{
大语言模型在医疗健康领域展现出广阔的应用前景，经过医学数据微调的模型在临床问答和决策支持等任务上已接近专业水平。然而，医疗决策直接关系患者生命安全，模型输出中的事实性错误、偏见或有害建议可能造成严重后果。现有安全评估研究主要针对通用场景或单体模型，在医疗专业性、多智能体协作场景覆盖以及真实临床数据验证等方面存在明显不足。

为解决上述问题，本文从三个方面开展研究。首先，提出面向医疗场景的大语言模型安全评估框架，该框架基于AutoGen和Ollama构建模块化的分层架构，包含模型集成层、智能体交互层和应用评估层，设计了六类功能智能体协同工作，支持数据投毒、后门攻击、提示词注入和越狱攻击四类威胁的自动化评估，并引入人机回环验证机制以保障评估可靠性。其次，基于通用医疗基准数据集，对五个医疗专用模型和两个通用基线模型开展单体对抗鲁棒性评估。实验结果表明，医疗专用模型的整体脆弱性高于通用模型，领域微调在提升专业能力的同时削弱了安全对齐强度；组合提示词攻击的平均攻击成功率达60.08\%，且三类攻击均具有高隐蔽性，基于文本表面质量的检测手段难以有效识别。最后，构建基于真实类风湿关节炎临床数据的五智能体级联架构，对级联错误传播假设进行实证检验。实验测得各位置的错误放大因子均大于1，系统级总放大因子在高强度攻击下达到8.4至9.1倍，验证了上游错误通过上下文锚定效应在协作链中逐步放大的机制。实验还发现链尾智能体的安全状态对系统整体风险具有决定性影响，保护该节点可将系统攻击成功率从最高95.8\%降至基线水平。人工专家复核实验的Cohen's Kappa系数为0.802，验证了自动化评估结果的可靠性。
}

\ckeywords{大语言模型；安全评估；多智能体系统；对抗性攻击；医疗人工智能}

\eabstract{This is abstract.

英文摘要字体为Times New Roman，小四，1.5倍行距。

英文摘要和关键词应与中文相对应。英语摘要用词应准确，使用本学科通用的词汇；摘要中主语（作用）常常省略，因而一般使用被动语态；应使用正确的时态，并要注意主、谓语的一致，必要的冠词不能省略。}

\ekeywords{Keyword1, Keyword2, Keyword3}
